% ============================================================================
% GEDES EXTENDED ABSTRACT — follows GEDESabstract_sampleFormat_en.docx:
%   A4, 4 pages, margins top/bottom 25 mm side 20 mm, Times 10 pt body,
%   single-spaced, fully justified, 1-letter paragraph indent;
%   TWO-COLUMN body (columnsep 21 pt) under a full-width title band
%   (matches the template's num=2 column section);
%   every figure AND the results table sit INSIDE one column (single-column
%   figure/table, no figure*/table*); the wide 4-probe timeline is scaled
%   down to column width (labels small but complete);
%   title centered 14 pt bold ALL CAPS; student block right-justified;
%   headings "1. Xxx" 10 pt bold flush left; figure captions "Fig. 1. ..."
%   centered BELOW the figure; numbered references [1], [2,3];
%   PDF rules: embed all fonts (pdflatex+newtx does), NO bookmarks
%   (hence: no hyperref), no security.
% All numbers from the certified 2026-07-06 set (code/results/*.json).
% American English throughout.
% OVERLEAF-READY / SELF-CONTAINED: every figure is a real file in ./figures/
%   (no symlinks, no ../ paths). To use: upload this .tex + the figures/ folder
%   to Overleaf (Times/newtx is on Overleaf by default). To refresh the figures
%   from the pipeline, re-copy from code/results/figures/ (see README note).
% Page limit: GEDES template = A4, 4 pages max. Cross-refs use \ref so the
%   figure numbers stay correct if you add/remove a float.
% Build: /Library/TeX/texbin/pdflatex gedes_abstract.tex  (x2)
% ============================================================================
\documentclass[10pt,a4paper,twocolumn]{article}
\usepackage[T1]{fontenc}
\usepackage[utf8]{inputenc}
\usepackage{newtxtext,newtxmath}          % Times; embeds Type-1 fonts
\usepackage[top=25mm,bottom=25mm,left=20mm,right=20mm]{geometry}
\usepackage{graphicx}
\usepackage{amsmath,bm}
\usepackage[labelformat=simple,labelsep=period,justification=centering,
            font=small]{caption}
\renewcommand{\figurename}{Fig.}
\usepackage{titlesec}
\usepackage{indentfirst}   % GEDES: ALL paragraphs indented, incl. first after a heading
\titleformat{\section}{\normalsize\bfseries}{\thesection.}{0.5em}{}
\titlespacing{\section}{0pt}{6pt}{3pt}
\titleformat{\subsection}{\normalsize\bfseries}{\thesubsection.}{0.5em}{}
\titlespacing{\subsection}{0pt}{4pt}{2pt}
\setlength{\parindent}{1em}
\setlength{\columnsep}{21pt}  % GEDES template: 420 twip
\pagestyle{empty}
% tighten float/caption spacing so the 4 figures + text hold 4 A4 pages
\usepackage{float}
\usepackage{stfloats}   % allow full-width figure*/table* at page BOTTOM too,
                        % so they pack instead of dumping onto their own pages
\usepackage{caption}
\captionsetup{skip=3pt}
\setlength{\floatsep}{6pt plus 2pt minus 2pt}
\setlength{\textfloatsep}{7pt plus 2pt minus 2pt}
\setlength{\intextsep}{6pt plus 2pt minus 2pt}
% Let a page carry more float area (default caps orphan big floats to their
% own page): raise the top/bottom/float-page fractions for both single- and
% double-column (starred) floats, and lower the minimum text fraction.
\renewcommand{\topfraction}{0.92}
\renewcommand{\bottomfraction}{0.75}
\renewcommand{\textfraction}{0.07}
\renewcommand{\floatpagefraction}{0.72}
\renewcommand{\dbltopfraction}{0.92}
\renewcommand{\dblfloatpagefraction}{0.72}
\setcounter{topnumber}{3}
\setcounter{bottomnumber}{2}
\setcounter{totalnumber}{5}
\setcounter{dbltopnumber}{3}

\begin{document}

% ── Full-width title band above the two columns (GEDES: title 1-col, body 2-col)
\twocolumn[{%
  \begin{center}
  {\fontsize{14}{17}\selectfont\bfseries
  DIFFERENCE OF LUNAR REGOLITH THERMAL CONDUCTIVITY $K_d$
  AT THE APOLLO 15 AND 17 HEAT-FLOW BOREHOLES}
  \end{center}
  \vspace{\baselineskip}
  \begin{flushright}
  Student No.: 24M58378\\
  Name: Ramon III P.\ GREGORIO\\
  Supervisor: Yasuko KASAI\\
  Co-supervisor: Richard LARSSON
  \end{flushright}
  \vspace{1.8\baselineskip}
}]

\section{Introduction}
Whether water ice can survive below the lunar surface, and how cold the
shallow subsurface stays, both depend on one quantity: the temperature
beneath the surface. No instrument measures that temperature directly,
so it has to be modeled. The number that controls such a model at depth is
the thermal conductivity of the regolith, the blanket of crushed rock
and dust that covers the Moon, written $K_d$. In everyday terms $K_d$
sets how well that blanket insulates the interior, so getting it wrong
shifts the whole subsurface temperature profile that mapping missions
depend on. Today every model uses a single value for the entire Moon,
$K_d = 3.4$ mW m$^{-1}$ K$^{-1}$, fit by Hayne et al.\ to orbital
measurements of how the surface cools at night [1]. But that
calibration only feels the top ${\sim}30$ cm; at the one-to-two-meter
depths that matter, the single global value is an educated
extrapolation, never checked against a subsurface measurement.

The subsurface has actually been measured at meter depths in only two
places: the Apollo 15 and 17 Heat Flow Experiments (HFE), temperature
probes that astronauts buried at each site, which recorded from 1971 to
1977 [2] and were recently restored to modern usability [3].
Figure~\ref{fig:sites} locates the two on the nearside, over the modeled
mean surface temperature and LOLA relief, with Apollo~15 at the
Hadley--Apennine mare/highland margin and Apollo~17 in the
Taurus--Littrow basaltic valley. This work addresses one deliberately
narrow question: what deep conductivity each of the two boreholes
requires when the standard model is held fixed in every other respect.
The answer feeds the temperature models that upcoming subsurface-mapping
missions depend on, in particular the TSUKIMI terahertz surveyor for
lunar water and resources [5], which cannot measure conductivity from
orbit and must assume it.

\begin{figure}[tbp]
\centering
\includegraphics[width=\linewidth]{figures/fig_context_map.pdf}
\caption{The two Apollo HFE sites on the nearside, over the modeled
mean surface temperature with LOLA shaded relief.}
\label{fig:sites}
\end{figure}

\section{Data and method}
\subsection{Data}
Reducing each sensor's multi-year record to one equilibrium
temperature is a step this study treats carefully, and it is where the
retrieval starts. For every sensor an automatic criterion selects the
\emph{stability window}, the longest late-mission stretch over which
the temperature trend has flattened below a strict $0.08$ K yr$^{-1}$
threshold, chosen after the drilling disturbance has decayed and clear
of documented operational events; the window mean is that sensor's
$T_{eq}$ (the colored bars in Fig.~\ref{fig:record}).
The vertical bands mark documented operational disturbances (drilling,
ALSEP outages, major station events) that are excluded from the window,
and the dashed line in each block fits the deepest sensor, whose small
residual warming (${+}0.01$ to ${+}0.08$ K yr$^{-1}$) is the long-term
subsurface trend analyzed by Nagihara et al.\ [3]. Sensors
shallower than 80 cm are excluded because the borestem conducts the
surface wave downward (gray rows); this leaves 7 deep sensors at
Apollo 15 and 16 at Apollo 17, spanning 80--234 cm. The basal boundary
condition is the revised Langseth et al.\ [2] borehole heat flux,
$Q_b = 21$ and 16 mW m$^{-2}$.

\begin{figure}[tbp]
\centering
\includegraphics[width=\linewidth]{figures/fig_apollo_timeline.pdf}
\caption{Restored 1971--1977 HFE record, all four probes. Per block:
every sensor's raw trace (depth-colored; borestem gray) with the
deepest-sensor drift fit dashed, above each sensor's retrieval
stability window (window-mean $T_{eq}$ at right). Vertical bands are
excluded disturbances; the 7 (A15) + 16 (A17) sub-borestem sensors are
the retrieval set.}
\label{fig:record}
\end{figure}

\subsection{Forward model and the flux-anchored solver}
The temperature column obeys the one-dimensional heat equation with the
Hayne conductivity form
\[
K(T,z) = [K_d - (K_d - K_s)e^{-z/H}]\,[1 + \chi (T/350\,\mathrm{K})^3],
\]
whose deep asymptote $K_d$ is the single free parameter [1]. The
observable is the periodic steady state, which brute-force
time-stepping reaches only after ${\sim}10^3$ simulated lunations per
solve. The methodological contribution of this
work is a \textit{flux-anchored} solver: in periodic steady state the
cycle-mean flux is constant in depth and pinned to $Q_b$, so the
two-variable heat equation collapses below the diurnal skin to a
one-variable closure ODE,
\begin{equation}
\frac{d\langle T\rangle}{dz}
   = \frac{Q_b - u_{\mathrm{rect}}(z)}{K(\langle T\rangle, z)},
\end{equation}
where $u_{\mathrm{rect}}$ is a small correction left by the day-night
swing. In plain terms, once the column settles into a steady repeating
cycle the average downward heat is the same at every depth, so the deep
temperature can be rebuilt from a single anchor point instead of
simulated step by step. The solver (Fig.~\ref{fig:method}) time-steps
only the thin skin, reads the cycle-mean temperature at an anchor depth
of 0.55 m, and reconstructs the deep column by integrating Eq.~(1). This cuts the
cost of one solve by a ${\sim}21\times$ algorithmic factor, or
${\sim}2500\times$ with the compiled kernel, at certified-equal
answers. Because these are ratios of two methods run on the same
machine, they do not depend on the hardware, and the speedup is what
makes the 1500-fold bootstrap of Fig.~\ref{fig:robust}a and the
point-by-point flux-plane map of Fig.~\ref{fig:robust}b affordable.

\begin{figure}[tbp]
\centering
\includegraphics[width=\linewidth]{figures/fig_method.pdf}
\caption{The flux-anchored solver. Step A time-steps only the diurnal
skin and reads the cycle-mean at the anchor; Step B reconstructs the
deep column from the closure ODE, Eq.~(1); the loop repeats until the
anchor stops drifting.}
\label{fig:method}
\end{figure}

\subsection{Retrieval and uncertainty}
For each site, $K_d$ is swept across a dense grid; each candidate is
solved to periodic equilibrium and scored by the root-mean-square
misfit to the deep sensors. The retrieved $K_d^{*}$ is the vertex of
the resulting misfit bowl, shown with its 95\% bootstrap interval as a
star and horizontal bar in Fig.~\ref{fig:sweep}. Uncertainty is
quantified by a non-parametric bootstrap: 1500 resamples of the sensor
set, with placement depths jittered by their 2.5 cm uncertainty [3],
each re-retrieved in full, an analysis made routine by the fast
solver. Systematic errors are budgeted separately, by re-running the
retrieval over the literature uncertainty range of each fixed input:
with the statistical term folded in, the per-site $1\sigma$ totals
are 1.88 (A15) and 3.88 (A17)
mW m$^{-1}$ K$^{-1}$, dominated by the basal flux ($\pm 1.6$ at both
sites) and, at Apollo 17, by the even larger surface-albedo term,
while the statistical and solver terms are minor ($\leq 0.53$ and 0.04).
The systematic terms are bounded literature ranges, not Gaussian
errors, so we keep them at $1\sigma$: tripling them to $3\sigma$ would
push $K_d$ outside its physical range (negative at Apollo~15) and is
not meaningful. Even so, the main result, both sites above the global
value, survives $3\sigma$ of the bootstrap; only the site-to-site
contrast stays marginal.

\begin{figure}[tbp]
\centering
\includegraphics[width=\linewidth]{figures/fig_kd_sweep.pdf}
\caption{Deep-sensor RMSE versus $K_d$. Stars mark the retrieved
$K_d^{*}$ (bowl vertex); horizontal bars are the $1\sigma$ statistical
(bootstrap) uncertainty, asymmetric where the bowl is flat; the dashed
vertical is the Hayne [1] global $3.4$, which both sites exceed.}
\label{fig:sweep}
\end{figure}

\section{Results}
Both boreholes favor deep conductivities well above the
Hayne et al.\ [1] global value, and markedly different from each
other (Fig.~\ref{fig:sweep}):
\begin{equation}
\begin{aligned}
K_d^{*}(\mathrm{A15}) &= 4.60^{+0.86}_{-0.21}\,[4.18,\,6.96],\\
K_d^{*}(\mathrm{A17}) &= 7.08^{+0.53}_{-0.52}\,[6.16,\,8.07]
\end{aligned}
\end{equation}
in mW m$^{-1}$ K$^{-1}$: the asymmetric $^{+}_{-}$ is the $1\sigma$
statistical (bootstrap) uncertainty, and the brackets are the 95\%
confidence interval. Propagating the published ranges of the fixed
inputs, dominated by the basal flux and, at Apollo 17, the surface
albedo, adds a systematic uncertainty that raises the total $1\sigma$
budget to $\pm1.9$ (A15) and $\pm3.9$ (A17). The shapes of
the two misfit bowls in Fig.~\ref{fig:sweep} explain the error shapes:
the Apollo 15 bowl is nearly flat toward high $K_d$, so its interval
stretches far upward ($+0.86$) but little downward ($-0.21$), whereas
the deeper, more symmetric Apollo 17 bowl yields a nearly symmetric
interval. The global
$K_d = 3.4$ lies outside both sites' statistical
intervals: it misfits the deep sensors by 1.09 K (A15) and 0.89 K
(A17), whereas the site fits reduce this to 1.00 and 0.40 K.
Figure~\ref{fig:profile} shows these fits directly against the sensor
temperatures, with the independent Mart\'{\i}nez \& Siegler [7]
$T,\rho$-dependent model tracking the Hayne global curve: the
deep-sensor gap is a property of the data, not of the Hayne
parameterization. That
exclusion is statistical, at the published inputs; the systematic
budget of Section~2.3 ($\pm1.9$ and $\pm3.9$ mW m$^{-1}$ K$^{-1}$)
weakens it to ${\sim}1\sigma$, so the model-robust result is the
consistent upward direction rather than a decisive rejection. The
inter-site contrast has bootstrap median $+2.31$ mW m$^{-1}$ K$^{-1}$
with 95\% interval $[-0.12, 3.56]$ and one-sided bootstrap tail
$P_{\rm boot}(\Delta K_d \le 0) = 0.031$: marginal at the 95\% level, and
reported here as such rather than overstated.
Table~\ref{tab:results} collects the retrieved values.

\begin{figure}[tbp]
\centering
\includegraphics[width=\linewidth]{figures/fig_profile_fit.pdf}
\caption{The fit. Annual-mean temperature vs depth: deep sensors (filled,
used) and shallow sensors (open, borestem-excluded), against the Hayne
[1] global-$3.4$ model and the Mart\'{\i}nez \& Siegler [7] $T,\rho$
model. Both global forms run cold of the deep sensors, which the per-site
$K_d$ closes.}
\label{fig:profile}
\end{figure}

\begin{table}[tbp]
\centering
\caption{Retrieved per-site $K_d^{*}$ and fit quality ($K_d$ in
mW\,m$^{-1}$\,K$^{-1}$; $Q_b$ in mW\,m$^{-2}$; RMSE in K). Brackets on
$K_d^{*}$ are the $95\%$ bootstrap confidence interval; the total
column is the $1\sigma$ total error budget (statistical plus
systematic, dominated by the basal flux and albedo). RMSE is against
the $\geq 80$ cm sensors, for the site fit and the global $3.4$.}
\label{tab:results}
\vspace{1pt}
{\footnotesize\setlength{\tabcolsep}{3.5pt}\renewcommand{\arraystretch}{1.45}
\begin{tabular}{lccccc}
\hline
Site & $Q_b$ & $K_d^{*}$ & total & RMSE & RMSE \\
     &       & [95\% CI] & ($1\sigma$) & fit  & 3.4  \\
\hline
A15 & 21 & $4.60\,[4.18, 6.96]$ & $\pm1.9$ & 1.00 & 1.09 \\
A17 & 16 & $7.08\,[6.16, 8.07]$ & $\pm3.9$ & 0.40 & 0.89 \\
\hline
\multicolumn{6}{l}{Contrast A17$-$A15: $+2.31$ $[-0.12, 3.56]$,
$P_{\rm boot}{=}0.031$}\\
\hline
\end{tabular}}
\end{table}

\section{Robustness of the contrast}
The dominant systematic is the basal flux itself, which is debated at
both sites [2,4]. Rather than assume the usual proportional
$K_d$--$Q_b$ degeneracy, this work maps it directly, re-retrieving
$K_d^{*}$ at every point of the two-site flux plane. The site responses are
opposite: $K_d^{*}$(A15) rises with its flux while $K_d^{*}$(A17)
falls, so the contrast never changes sign. Figure~\ref{fig:robust}b
shows it as a heatmap of the contrast over the two basal fluxes: warm
(Apollo~17 more conductive) across the whole plausible plane, with the
``no-difference'' line ($\Delta = 0$) reaching only the far
low-A15/high-A17 corner. The dashed contours grade the significance in
units of the bootstrap standard error: at the Langseth et al.\ [2]
fluxes (dot; 21 and 16 mW m$^{-2}$) the contrast is $+2.5$~mW m$^{-1}$ K$^{-1}$, and it only grows away from
that point. The square marks
the independent reanalysis of Saito et al.\ [4], which lowers the
Apollo 15 flux by ${\sim}30\%$: there the contrast nearly doubles,
to $+4.8$. The shaded teal region is the Bayesian MCMC posterior for
the two fluxes (68\% and 95\% credible); it keeps the ordering
$P[K_d^{*}(\mathrm{A17}) > K_d^{*}(\mathrm{A15})] = 99.2$--$99.4\%$ for
prior widths from $\pm 10\%$ to $\pm 40\%$, and its medians (star; 17.8 and
13.9 mW m$^{-2}$) sit slightly below the Langseth et al.\ [2] values,
in the direction Saito et al.\ argued.

\begin{figure}[tbp]
\centering
\includegraphics[width=\linewidth]{figures/fig_robustness.pdf}
\caption{Per-site \emph{statistical} uncertainty (a, sensor-resampling
bootstrap of $K_d^{*}$; the larger systematic uncertainty from the basal
flux and albedo is not shown here) and robustness of the site contrast
to the basal flux (b).}
\label{fig:robust}
\end{figure}

A factor-${\sim}1.5$ contrast between the two sites is geologically
unremarkable: Apollo core porosities alone span an effective-medium
conductivity factor of ${\sim}1.5$. The setting is visible in
Fig.~\ref{fig:sites}: in the albedo and topography layers Apollo~17
sits inside the dark, low-lying basaltic embayment of Taurus--Littrow,
while Apollo~15 straddles the bright mare/highland margin at
Hadley--Apennine. The two sites have different regolith column
histories, so a difference in their meter-scale conductivities is
expected rather than anomalous.

\section{Conclusions}
(1) The Apollo boreholes favor site-resolved deep conductivities,
$4.60$ and $7.08$ mW m$^{-1}$ K$^{-1}$, each placing the global 3.4
outside its bootstrap 95\% interval at the published inputs
(${\sim}1\sigma$ once systematics are folded in), a consistent
upward departure that argues against a single Moon-wide $K_d$ for
meter-scale thermal models. (2) The inter-site contrast is marginal in
magnitude (one-sided $P_{\rm boot}=0.031$) but directionally robust
across the entire plausible basal-flux domain (MCMC ordering
probability 99.2--99.4\%).
(3) The flux-anchored solver makes exhaustive uncertainty analysis of
borehole-constrained thermal models routine (${\sim}2500\times$ faster
per solve than brute force, certified to reproduce the brute-force
answer). (4) The retrievals are conditional
on the fixed Hayne conductivity form, and the \emph{absolute} values
scale with its radiative coefficient $\chi = 2.7$: re-retrieving at the
alternative Diviner normalization $\chi = 1.48$ [6] drives both misfit
minima outside the physical sweep window, so the absolute $K_d$ is
unconstrained there, yet Apollo 17 still demands the higher conductivity,
so the \emph{direction} of the contrast is insensitive to the radiative
normalization. Relaxing the mixing-depth parameter $H$ likewise only
widens the Apollo 17 uncertainty (to ${\pm}3.5$ mW m$^{-1}$ K$^{-1}$)
without reversing the ordering. These are caveats this work states
rather than hides. (5) The retrieved per-site temperature columns supply
a meter-scale temperature constraint for the forward models of
TSUKIMI-class subsurface sounders, which cannot measure $K_d$ from
orbit [5].

Taken together, the two boreholes show that the Moon's insulating
blanket is not the same everywhere: a single global conductivity
misjudges the subsurface temperature at the Apollo depths by about
1~K, the same size as the signal the next generation of sounders must
read. The practical payoff is in the method's speed. Because one solve
is now about $2500\times$ cheaper, the same retrieval can be run not
only at two boreholes but pixel by pixel across the Moon, turning
scattered ground truth into maps of subsurface temperature, and from
those, maps of where the ground stays cold enough for water ice to
persist. That is the direction this work opens: from two points to a
global picture, with the deep conductivity measured rather than
assumed.

\section*{References}
\begin{list}{}{\setlength{\leftmargin}{1.5em}\setlength{\itemindent}{-1.5em}
\setlength{\itemsep}{0pt}\setlength{\parsep}{0pt}}
\item[] [1] Hayne, P. O., Bandfield, J. L., Siegler, M. A., et al.,
Global regolith thermophysical properties of the Moon from the Diviner
Lunar Radiometer Experiment, J. Geophys. Res. Planets, vol.122, no.12,
2017, pp. 2371--2400.
\item[] [2] Langseth, M. G., Keihm, S. J. and Peters, K., Revised lunar
heat-flow values, Proc. Lunar Sci. Conf. 7, 1976, pp. 3143--3171.
\item[] [3] Nagihara, S., Kiefer, W. S., Taylor, P. T., Williams, D. R.
and Nakamura, Y., Examination of the long-term subsurface warming at the
Apollo 15 and 17 sites from the restored heat-flow data, J. Geophys.
Res. Planets, vol.123, no.5, 2018, pp. 1125--1139.
\item[] [4] Saito, Y., Tanaka, S., Takita, J., Horai, K. and
Hagermann, A., Lunar heat flow experiment for SELENE-2, Adv. Space Res.,
vol.40, 2007, pp. 1601--1607.
\item[] [5] Kasai, Y., Sato, T., Yamada, T., Hasegawa, Y. and the
TSUKIMI team, TSUKIMI: Lunar Terahertz Surveyor for Kilometer-scale
Mapping, 44th COSPAR Sci. Assembly, 2022, p. 291.
\item[] [6] Vasavada, A. R., Bandfield, J. L., Greenhagen, B. T., et al.,
Lunar equatorial surface temperatures and regolith properties from the
Diviner Lunar Radiometer Experiment, J. Geophys. Res. Planets, vol.117,
2012, E00H18.
\item[] [7] Mart\'{\i}nez, A. and Siegler, M. A., A global thermal
conductivity model for the surface of the Moon, J. Geophys. Res.
Planets, vol.126, no.11, 2021, e2021JE006829.
\end{list}

\end{document}
